% Copyright (c) 2026 Eugene Joseph M. Ragasa.
% SPDX-License-Identifier: Apache-2.0
% Author: Eugene Joseph M. Ragasa
% Formal proof completed: 2026-10-19
\documentclass[11pt]{article}
\usepackage[T1]{fontenc}\usepackage{amsmath,amssymb}\usepackage[margin=1in]{geometry}\usepackage[hidelinks]{hyperref}
\title{Proof Dossier: Uniform Enclosure Ellipse Separation}\author{Eugene Joseph M. Ragasa}\date{Formal proof completed: 2026-10-19}
\begin{document}\maketitle
\section{Identity and status}
ID \texttt{BG-CERT-THM-008}; status \texttt{proved}. Lean declaration:
\par\noindent{\tiny\path{ProjectKoios.BandGauge.Certificates.}}\\
{\tiny\path{disjoint_fullEllipses_of_enclosureCheck_eq_true}}.
\section{Statement and proof}
For every real center in the checked rational interval, the exact spectral and operator
full ellipses are disjoint. The upper endpoint gives both the largest operator radius and
the smallest center gap. Checked nonnegativity converts the strict squared rational
inequality to the reusable strict-separation hypothesis.
\section{Boundary}
This supports incompatibility near Paper A manuscript line 857 at owner commit
\begin{center}\small\texttt{99a46cb59d0c6f27ada7d709f194996ef9619d9f}.\end{center}
No concrete field binding is claimed. The kernel axiom report is:
\begin{itemize}
\item \texttt{propext};
\item \texttt{Classical.choice}; and
\item \texttt{Quot.sound}.
\end{itemize}
There are no admissions. The proof is AI-assisted and kernel-checked.
\end{document}
